\section{The nilpotent deck group and small-trace projections}\label{nil:section}

We construct a torsion-free nilpotent deck group that serves both parts
of the argument. Its voltages must detect the short holonomy words
required by the geometric construction, while its central Fourier fibers
must carry nonzero projections of small weighted identity trace. We first
fix the word detector, then build the group and the projection field.
Section~\ref{sam:section} chooses the actual edge voltages and matrix
weights that meet the geometric and operator requirements simultaneously.

The constants \(C_*\) and \(C_\bullet\) in the geometric construction are
fixed throughout this section. Every constant below is independent of the
layer index and, unless explicitly a function of \(d\), independent of
\(\ell\). Write \(\ee(v)=\exp(2\pi\mathrm i v)\).
All discrete Hilbert spaces have counting measure. Traces on auxiliary
matrix indices use their ordinary sum, without division by their number.

\subsection{A uniform detector for short walks}

\begin{lemma}[Rank bound]\label{nil:rank}
Let \(S\) be a finite connected metric graph with positive edge lengths,
total length at most \(Cg\), and girth greater than \(g>0\).
Girth means the infimum of lengths of embedded circles, with value
\(+\infty\) for a tree. Put
\[
 B(C)=\max\{1,\lceil20C\rceil\}^{\,2}.
\]
Then \(b_1(S)\le B(C)\). Vertices of degree one and parallel edges are
allowed.
\end{lemma}

\begin{proof}
Give \(S\) its intrinsic path metric and put \(a=g/10\).
Choose a maximal \(a\)-separated set \(\mathcal P\), which is an \(a\)-net.
If it has at least two elements, each \(p\in\mathcal P\) is joined to
another by a shortest path of length at least \(a\), so its open ball
of radius \(a/2\) contains an interval of length \(a/2\). Those balls
are disjoint. Therefore
\[
 |\mathcal P|a/2\le \operatorname{length}(S)\le Cg,
 \qquad |\mathcal P|\le\max\{1,\lceil20C\rceil\}.
\]
The second bound also covers the singleton case.

Form an abstract graph \(K\) on \(\mathcal P\), with one edge for each
distinct pair at distance at most \(3a\), and map each edge to a
chosen shortest path in \(S\). This graph is connected and its map
surjects on fundamental groups, using a connector for basepoints.
Indeed subdivide a path of \(S\) into segments of length at most \(a\),
choose net points within distance \(a\) of their endpoints, and replace
each segment by the chosen path between its two net points. Each
replacement exists. The original segment, its two endpoint connectors,
and its replacement form a closed path of length at most \(6a<g\).
Such a path is null-homotopic: after subdivision and cancellation
of backtracks, a nonempty reduced closed path in a graph contains
an embedded circle no longer than itself. Thus replacement preserves
homotopy relative to the connected endpoints. Applying this to paths
and then loops proves the two assertions about \(K\), and hence
\[
 b_1(S)\le b_1(K)\le|\mathcal P|^2\le B(C).
\]
\end{proof}

\begin{lemma}[Bounded pivot words and a matrix detector]\label{nil:word}
Let \(W\) be a nonempty cyclically nonbacktracking closed walk in a
unit-edge graph of girth greater than \(g\), with length at most \(Cg\).
There is a spanning tree of its support such that, after its edge
variables are set equal to the identity, the holonomy word has a
nontrivial free reduction of length at most
\[
 L(C)=\lceil3CB(C)\rceil.
\]
If \(m>L(C)\), that word is not an identity on the group \(U_{\mathbb Z}\)
of upper unitriangular \(m\)-by-\(m\) integer matrices.
Every entry of any word in generic upper unitriangular matrices and
their inverses is a polynomial of total degree at most \(m-1\) in their
upper entries. Nonidentity is preserved if each edge variable is
multiplied on either side by a fixed upper unitriangular matrix.
\end{lemma}

\begin{proof}
Let \(S\) be the support and \(b=b_1(S)\). A degree-one vertex would
force a backtrack, including at the cyclic join, so \(S\) has minimum
degree at least two and \(b\ge1\). If all degrees are two, regard the
circle as one arc with a chosen endpoint. Otherwise suppress
degree-two vertices. Euler's formula and minimum degree three give
at most \(2b-2\) vertices and \(3b-3\) edges. In either case there
are at most \(3b\) arcs.

Arcs shorter than \(g/(3b)\) form a forest, since a circuit made of
them would have length less than \(g\). Extend this forest to a
spanning tree of the suppressed graph. From each complementary arc
delete one original edge, called its pivot. The remaining original
graph is a spanning tree \(T\). Cyclically start \(W\) at a retained
vertex. Nonbacktracking forces it to traverse each arc it enters
completely. Every pivot crossing belongs to a distinct traversal
of an arc of length at least \(g/(3b)\). There are consequently
at most \(3bC\le L(C)\) pivot crossings, by Lemma~\ref{nil:rank}.
Collapsing \(T\) is a homotopy equivalence to a bouquet of circles.
Its image word is nontrivial since a nonempty cyclically reduced
path in a graph is not null-homotopic. Free reduction cannot
increase the length.

For the matrix assertion we use a finite matrix truncation of the Magnus
power-series expansion \cite{Magnus1935}. Write a nontrivial reduced word
in maximal syllables:
\[
 w=x_{j_1}^{a_1}\cdots x_{j_k}^{a_k},\qquad
 a_s\in\mathbb Z\setminus\{0\},\quad j_s\ne j_{s+1}.
\]
In formal noncommutative power series set \(x_j=1+Z_j\), using
\((1+Z)^a=\sum_{h\ge0}\binom ah Z^h\) also for negative \(a\).
The coefficient of \(Z_{j_1}\cdots Z_{j_k}\) is
\(\prod_s a_s\ne0\). A contributing syllable supplies at most one
letter, since the target has no equal adjacent letters; thus all
\(k\) syllables must supply exactly one. For \(m>k\), put
\[
 Z_j=\sum_{\{s:j_s=j\}}E_{s,s+1}\in M_m(\mathbb Z).
\]
The \((1,k+1)\) entry detects exactly that monomial. It follows that
\(w(1+Z_j)\ne1\). Here \(k\le\operatorname{length}(w)<m\).

The inverse of \(1+X\), for strictly upper triangular \(X\), is
\(\sum_{h=0}^{m-1}(-X)^h\). Every contributing product in any
matrix entry follows strictly increasing matrix indices, and hence
uses at most \(m-1\) upper-entry variables regardless of word length.
This proves the degree bound. Multiplication by a fixed
unitriangular matrix is a polynomial coordinate bijection with
polynomial inverse. Separate such changes in the edge variables
preserve a nonzero polynomial.
\end{proof}

Fix once and for all \(m\ge2\) with
\begin{equation}\label{nil:choice-m}
 m>L(201C_*).
\end{equation}
Indeed a walk based in layer \(i\) of length at most
\(C_*H_i/\log_2d\) stays in layers \(j\) with
\(H_j\ge H_i-C_*H_i/\log_2d\). For sufficiently large \(\ell\),
the tower girth property therefore makes every circuit in its
support longer than \(H_i/(201\log_2d)\). Lemma~\ref{nil:word},
with this \(g\) and \(C=201C_*\), applies. Lemma~\ref{nil:rank}
can later be applied with any other fixed \(C\), without changing \(m\).

\subsection{An integer affine group and its central product}

The short-word test only requires a quotient onto \(U_{\mathbb Z}\),
so we can add coordinates without losing that detector. The affine
extension below is chosen for a second purpose: it admits a twisted
representation in which Schwartz coefficient functions correspond to
Schwartz integral kernels. A rank-one kernel then gives a continuous
convolution projection, which we will sample on integer coordinates.

Let \(\mathfrak n\) be the strictly upper triangular real
\(m\)-by-\(m\) matrices, let \(n_{\mathrm f}=m(m-1)/2\), and let
\(U=1+\mathfrak n\). Write \(x_{ij}\) for the coordinates on
\(\mathfrak n\) and \(\lambda_{ij}\) for their dual coordinates.
For \(b\in U\), put
\[
 A_b(x)=b(1+x)-1=(b-1)+bx,\qquad L_b(x)=bx.
\]
The maps \(A_b\) form an affine action: \(A_bA_{b'}=A_{bb'}\).
Let \(E\) be the semidirect product of additive affine functions
on \(\mathfrak n\) by \(U\), acting by \(f\mapsto f\circ A_{b^{-1}}\).
Use coordinates \((\lambda,b,z)\), with affine function
\(x\mapsto z+\lambda(x)\). Explicitly,
\begin{align}
 (\lambda,b,z)(\mu,b',w)
 &=\bigl(\lambda+\mu\circ L_{b^{-1}},bb',
             z+w+\mu(b^{-1}-1)\bigr),\label{nil:affine-law}\\
 (\lambda,b,z)^{-1}
 &=\bigl(-\lambda\circ L_b,b^{-1},
             -z-\lambda(b-1)\bigr).\label{nil:affine-inverse}
\end{align}
Constants \((0,1,z)\) are central. Their quotient \(\bar E\) has
coordinates \(\gamma=(\lambda,b)\) and section
\(\sigma(\gamma)=(\lambda,b,0)\). For \(\eta=(\mu,b')\), put
\[
 \gamma\eta=(\lambda+\mu\circ L_{b^{-1}},bb'),\qquad
 c(\gamma,\eta)=\mu(b^{-1}-1).
\]
Then
\begin{equation}\label{nil:cocycle}
 \sigma(\gamma)\sigma(\eta)
       =(0,1,c(\gamma,\eta))\sigma(\gamma\eta),
\end{equation}
and associativity gives
\[
 c(\gamma,\eta)+c(\gamma\eta,\zeta)
   =c(\eta,\zeta)+c(\gamma,\eta\zeta).
\]
The cocycle is normalized at the identity, and
\(c(\gamma,\gamma^{-1})=c(\gamma^{-1},\gamma)\).

\begin{lemma}[Grading, Haar measure, and integer coordinates]\label{nil:grading}
The laws above are polynomial with integer coefficients in
\((\lambda,b-1,z)\). Give the coordinates weights
\[
 \operatorname{wt}(b-1)_{ij}=j-i,\quad
 \operatorname{wt}\lambda_{ij}=m-(j-i),\quad
 \operatorname{wt}z=m.
\]
Scaling weight \(a\) by \(y^a\) defines an automorphism \(\Delta_y\),
for \(y>0\), and \(c(\Delta_y\gamma,\Delta_y\eta)=y^m c(\gamma,\eta)\).
Lebesgue coordinate measure is both left and right Haar measure on
\(E\) and \(\bar E\); inversion preserves it. The Jacobian of
\(\Delta_y\) on \(\bar E\) is \(y^{mn_{\mathrm f}}\).
The integer coordinate sets \(E_{\mathbb Z}\) and \(\bar E_{\mathbb Z}\)
are countable nilpotent torsion-free groups.
\end{lemma}

\begin{proof}
The finite geometric series for \(b^{-1}\) proves polynomiality
and integrality in Equations~\ref{nil:affine-law} and
\ref{nil:affine-inverse}. The usual superdiagonal dilation is
an automorphism of \(U\). On affine functions use
\(f(x)\mapsto y^mf(\Delta_{y^{-1}}x)\).
The identity \(\Delta_y A_b=A_{\Delta_y b}\Delta_y\) shows that
these maps intertwine the actions, giving the stated automorphism
and cocycle homogeneity. Each paired quotient coordinate
\((\lambda_{ij},(b-1)_{ij})\) has total weight \(m\), proving
the dilation Jacobian.

Order the entries of \(b-1\) by superdiagonal. Left and right
multiplication on \(U\) are triangular coordinate transformations
with diagonal entries one, and thus determinant one. The map
\(L_b\) on \(\mathfrak n\) has determinant one: on each column it
uses a unitriangular initial submatrix of \(b\).
For a fixed left factor in Equation~\ref{nil:affine-law}, the
\(\mu\)-coordinates are transformed by \(L_{b^{-1}}^*\), the
\(b'\)-coordinates by left multiplication, and the \(w\)-coordinate
has coefficient one. Cross terms do not alter the determinant.
For a fixed right factor, the \(\lambda,z\) coordinates are
translated by functions of \(b\), and \(b\) is right multiplied.
This gives determinant one in both cases, also on omitting \(z\).
Inversion exchanges these equally normalized left and right Haar
measures, or has absolute Jacobian one by the same triangular
calculation.

In a coordinate of weight \(a\), multiplication is additive plus
a polynomial in smaller-weight coordinates. Each correction
monomial uses both factors, by homogeneity and the identity laws.
The subgroups whose coordinates of weight less than \(a\) vanish
form a filtration \(E_{\ge a}\) with
\([E,E_{\ge a}]\subset E_{\ge a+1}\). Indeed multiplication
modulo weight \(a+1\) by an element of \(E_{\ge a}\) is simply
addition of its weight-\(a\) coordinates, so is central there.
All weights are at most \(m\), proving nilpotence, also in the
quotient and the integer groups. A positive power of an element
with first nonzero weight \(a\) multiplies its weight-\(a\) vector
by that positive integer. It cannot become zero. This proves
torsion-freeness.
\end{proof}

Take copies \(X_{\mathrm f},Y_{\mathrm f}\) of \(E_{\mathbb Z}\);
these are factor names, unrelated to the graphs. Define
\begin{equation}\label{nil:deck}
 F=(E_{\mathbb Z}\times E_{\mathbb Z})/
       \{((0,1,k),(0,1,-k)):k\in\mathbb Z\}.
\end{equation}
Equivalently its coordinates are
\((\gamma_X,\gamma_Y,z)\in\bar E_{\mathbb Z}^2\times\mathbb Z\),
with law
\[
 (\gamma_X,\gamma_Y,z)(\eta_X,\eta_Y,w)
 =(\gamma_X\eta_X,\gamma_Y\eta_Y,
       z+w+c(\gamma_X,\eta_X)+c(\gamma_Y,\eta_Y)).
\]
The coordinate map from the product sums its two central integers
and has precisely the kernel in Equation~\ref{nil:deck}.
The opposite-pair subgroup is primitive in \(\mathbb Z^2\), so
the surviving coordinate is an integer coordinate with no finite
quotient. The inherited grading and first-nonzero-weight proof
show directly that \(F\) is countable, nilpotent, and torsion-free.
Projection to the upper unitriangular coordinate of either factor
is a surjective group homomorphism to \(U_{\mathbb Z}\).
Write \(z_0=(1,1,1)\) for the common central generator and
\(\sigma_X,\sigma_Y\) for the two section maps into \(F\).

\subsection{Central Fourier transform and equivariant trace}

Use right translations \((\rho(g)\xi)(f)=\xi(fg)\).
The inversion unitary \((\mathcal I\xi)(f)=\xi(f^{-1})\) satisfies
\(\mathcal I\rho(g)\mathcal I=\lambda_F(g)\).
After inversion, Fourier transform the central coordinate by
\(\delta_j\mapsto(t\mapsto\ee(tj))\), with probability Lebesgue
measure on \(\mathbb T=\mathbb R/\mathbb Z\). This identifies
\begin{equation}\label{nil:fourier}
 \ell^2(F)\cong L^2\bigl(\mathbb T;
               \ell^2(\bar E_{\mathbb Z})\otimes
               \ell^2(\bar E_{\mathbb Z})\bigr).
\end{equation}
The central generator acts by \(\ee(t)\). In either factor, the
section acts by
\begin{equation}\label{nil:twisted-left}
 u_t(\gamma)\delta_\eta
       =\ee(t c(\gamma,\eta))\delta_{\gamma\eta}.
\end{equation}
Thus
\[
 u_t(\gamma)u_t(\eta)=\ee(t c(\gamma,\eta))u_t(\gamma\eta),
 \quad
 u_t(\gamma)^*=\ee(-t c(\gamma^{-1},\gamma))u_t(\gamma^{-1}).
\]
The two factor representations commute; the full quotient
cocycle is their sum. There are commuting right operators
\begin{equation}\label{nil:twisted-right}
 v_t(k)\delta_\eta=\ee(t c(\eta,k))\delta_{\eta k}.
\end{equation}
The cocycle identity proves that these commute with \(u_t\).
Under the Fourier transform of the central coordinate on the already
inverted copy of \(\ell^2(F)\), these are the images of right multiplication
by \(\sigma(k)\) on basis indices, namely \(\rho(\sigma(k)^{-1})\).
They permute quotient basis vectors up to unit phases.
Equivariance below means intertwining these right actions,
with the identity action on auxiliary indices. Untwisted left
or right regular equivariance obeys the same arguments.

\begin{lemma}[Equivariant trace and size]\label{nil:trace}
Let \(L\) be a countable group with a right action as above and let
\(I,J\) be finite or countable auxiliary sets. On
\(\mathcal H_I=\ell^2(L)\otimes\ell^2(I)\), define, for equivariant
positive bounded \(B\),
\[
 \operatorname{Tr}_I(B)
   =\sum_{i\in I}\langle B\delta_{(1,i)},\delta_{(1,i)}\rangle
       \in[0,\infty].
\]
For a bounded equivariant \(A:\mathcal H_I\to\mathcal H_J\),
\begin{equation}\label{nil:adjoint-square}
 \operatorname{Tr}_I(A^*A)=\operatorname{Tr}_J(AA^*).
\end{equation}
This trace is additive on positive operators, monotone, and
faithful on projections, and does not depend on the chosen
orthonormal basis of the auxiliary Hilbert space.
Equivariantly equivalent projections
have equal size. A bounded equivariant map injective on the
range of a projection \(P\), with image there contained in the
range of \(Q\), gives \(\operatorname{Tr}(P)\le\operatorname{Tr}(Q)\).
These statements apply to countable orthogonal sums and to each
Fourier fiber. Integrating a fiber trace gives the identity trace
before Fourier transformation.
\end{lemma}

\begin{proof}
For a matrix entry of \(A\), translation on the right by \(g^{-1}\)
gives
\[
 |A((g,j),(1,i))|=|A((1,j),(g^{-1},i))|.
\]
The phases have absolute value one and the auxiliary indices
are fixed. Tonelli's Theorem for countable nonnegative sums gives
\begin{align*}
 \operatorname{Tr}_I(A^*A)
 &=\sum_{i,j,g}|A((g,j),(1,i))|^2\\
 &=\sum_{j,i,g}|A((1,j),(g^{-1},i))|^2
  =\operatorname{Tr}_J(AA^*),
\end{align*}
including when the sums are infinite.
Additivity and monotonicity follow from positive diagonals.
Changing the auxiliary orthonormal basis is conjugation by an
equivariant unitary \(U\). Applying the adjoint-square identity
to \(UB^{1/2}\) gives
\(\operatorname{Tr}(UBU^*)=\operatorname{Tr}(B)\), proving
basis independence.
If a projection \(P\) has zero trace, then
\(\|P\delta_{(1,i)}\|^2=0\) for each \(i\).
Right translates span the Hilbert space, so \(P=0\).

The polar partial isometry of an equivariant bounded \(A\) is
equivariant: it is the strong limit of
\(A(|A|+1/n)^{-1}\). Its initial and final projections are those
of \((\ker A)^\perp\) and \(\overline{\operatorname{ran}A}\).
Apply Equation~\ref{nil:adjoint-square} to it. This proves
projection equivalence and, on applying it to \(AP\), the
comparison, because its initial projection is \(P\) and its
final projection is at most \(Q\). Countable orthogonal sums
are handled by summing identity diagonals.

A vector at central coordinate zero becomes a constant vector
field in Equation~\ref{nil:fourier}. Its diagonal is the integral
of its fiber diagonal. Sum nonnegative diagonals and use Tonelli
for the integrated assertion, including countably many
auxiliary indices.
\end{proof}

A measurable uniformly bounded field on a separable Hilbert
space defines a bounded operator on its circle \(L^2\) space
with norm the essential supremum of the fiber norms.
Integration proves the upper inequality; for the lower one,
use a countable dense set of unit vectors restricted to
measurable circle sets where their images exceed a given norm.
Products and adjoints act pointwise. Uniform polynomial
approximation gives pointwise continuous functional calculus
for uniformly bounded self-adjoint fields. Thus a fixed
separating continuous function or resolvent contour gives the
same projection before and after decomposition.

\subsection{A Schwartz convolution projection}

Schwartz matrix coefficients and rank-one operators for Heisenberg
representations are treated by Rieffel
\cite[Lemma~2.3, Corollary~2.4, and Theorem~2.18]{Rieffel1988}.
Here we derive the coefficient-to-kernel correspondence for the affine
group above, whose quotient \(\bar E\) is generally nonabelian.

Equip \(\bar E\) with the Haar measure of Lemma~\ref{nil:grading}.
Schwartz functions refer to its Euclidean coordinates
\((\lambda,b-1)\), or to the specified product coordinates.
For \(\epsilon\in\{1,-1\}\), define
\begin{align*}
 (J*_\epsilon K)(h)
  &=\int_{\bar E}J(a)K(a^{-1}h)
          \ee\bigl(\epsilon c(a,a^{-1}h)\bigr)\,da,\\
 J^{*_\epsilon}(h)
  &=\overline{J(h^{-1})}
          \ee\bigl(-\epsilon c(h^{-1},h)\bigr).
\end{align*}
The cocycle identity gives associativity and the involution
identities. The following kernel realization verifies the
operations and gives the idempotents we need.

\begin{lemma}[Schwartz idempotents]\label{nil:schwartz}
For each sign \(\epsilon\), there is a Schwartz function
\(J_\epsilon\) on \(\bar E\) with
\[
 J_\epsilon^{*_\epsilon}=J_\epsilon,\qquad
 J_\epsilon *_\epsilon J_\epsilon=J_\epsilon,\qquad
 \|J_\epsilon\|_{L^2(\bar E)}=1.
\]
\end{lemma}

\begin{proof}
For \(\gamma=(\lambda,b)\), define on \(L^2(\mathfrak n)\)
\[
 (\pi_\epsilon(\gamma)\phi)(x)
       =\ee(\epsilon\lambda(x))\phi(A_{b^{-1}}x).
\]
The Jacobian of \(A_b\) is one, so these operators are unitary.
Direct substitution in Equation~\ref{nil:affine-law} gives
\[
 \pi_\epsilon(\gamma)\pi_\epsilon(\eta)
       =\ee(\epsilon c(\gamma,\eta))\pi_\epsilon(\gamma\eta).
\]
For a Schwartz \(J\), its integrated operator is bounded by
\(\|J\|_1\). Denote partial Fourier transform with positive sign by
\[
 \widehat J_+(\xi,b)=\int_{\mathfrak n^*}
                         J(\lambda,b)\ee(\lambda(\xi))\,d\lambda.
\]
Changing variables
\[
 y=b^{-1}(1+x)-1,\qquad b=(1+x)(1+y)^{-1}
\]
in that integrated operator gives its kernel
\begin{equation}\label{nil:kernel}
 \mathcal K_J(x,y)=
   \widehat J_+\bigl(\epsilon x,(1+x)(1+y)^{-1}\bigr).
\end{equation}
For fixed \(x\), the change from \(b-1\) to \(y\) has absolute
Jacobian one: it is inversion followed by multiplication in \(U\),
in its Lebesgue Haar coordinates.

Partial Fourier transform is a bijection on Schwartz spaces
and an \(L^2\)-isometry in the displayed normalization
\cite[Chapter 8]{Folland1999}.
The coordinate map
\[
 (x,y)\longmapsto
       \bigl(\epsilon x,(1+x)(1+y)^{-1}-1\bigr)
\]
and its inverse are polynomial. Pullback by such a map preserves
Schwartz functions bijectively: derivatives introduce polynomial
factors, and its polynomial inverse bounds any polynomial weight
in the source by a polynomial weight in the target. Its absolute
Jacobian is one as just shown. Consequently
\(J\mapsto\mathcal K_J\) is a bijection from Schwartz coefficients
to Schwartz kernels and an \(L^2\)-isometry.
It is also injective as an operator realization. If a continuous
Schwartz kernel gives the zero operator, its pairings against
all products of compactly supported smooth test functions vanish,
and approximations to point masses show the kernel vanishes.

Twisted convolution and involution preserve Schwartz coefficients.
Indeed
\[
 (a,h)\longmapsto J(a)K(a^{-1}h)
          \ee\bigl(\epsilon c(a,a^{-1}h)\bigr)
\]
is jointly Schwartz. The map \((a,h)\mapsto(a,a^{-1}h)\) is a
polynomial bijection with polynomial inverse, and differentiation
of the phase introduces only polynomial factors. Integration in
\(a\) preserves rapid decay of all derivatives in \(h\).
The same reasoning with inversion proves the involution assertion.
Absolute integrability permits Fubini, so the integrated
representations carry these operations to products and adjoints.

Choose \(\phi(x)=2^{n_{\mathrm f}/4}\exp(-\pi|x|^2)\), of \(L^2\)
norm one. The orthogonal projection onto its span has Schwartz
kernel \(\phi(x)\overline{\phi(y)}\). Let \(J_\epsilon\) be its
unique coefficient under Equation~\ref{nil:kernel}.
Product and adjoint correspondence and injectivity give the two
identities. The \(L^2\)-isometry gives
\[
 \|J_\epsilon\|_2^2
   =\int|\phi(x)|^2|\phi(y)|^2\,dx\,dy=1.
\]
\end{proof}

\subsection{The lattice limit in both coefficient norms}

Put \(L=\bar E_{\mathbb Z}\). For real \(\theta\), define
\(C_r^*(L,\theta c)\) to be the norm closure on \(\ell^2(L)\)
of the span of the operators in Equation~\ref{nil:twisted-left}
with \(t=\theta\). For \(a\in\ell^1(L)\), the series
\(\sum_\gamma a_\gamma u_\theta(\gamma)\) converges in norm
with norm at most \(\sum_\gamma|a_\gamma|\).
Its coefficient \(\ell^2\) norm equals its norm on \(\delta_1\),
because \(c(\gamma,1)=0\).
The identity diagonal is its identity coefficient and is the
trace of Lemma~\ref{nil:trace}. In particular
\(\operatorname{tr}_\theta(A^*A)=\|A\delta_1\|^2\).

The two coefficient norms have different roles in passing from the
continuous idempotent to the lattice. The \(\ell^1\) defect controls
the operator-norm correction by functional calculus. The \(\ell^2\)
defect, compared with the sampled coefficient norm, prevents the
resulting projection from being zero, while the squared \(\ell^2\)
size bounds its identity trace. The next lemma supplies both controls.

\begin{lemma}[Weighted rectangular Riemann sums]\label{nil:riemann}
For positive integer weights \(w_1,\ldots,w_k\), set
\[
 \mathcal L_y=\prod_{a=1}^k y^{w_a}\mathbb Z,\qquad
 v_y=y^{w_1+\cdots+w_k},\qquad 0<y\le1.
\]
For Schwartz \(f\) on \(\mathbb R^k\),
\[
 v_y\sum_{a\in\mathcal L_y}f(a)\longrightarrow
                   \int_{\mathbb R^k}f(x)\,dx.
\]
The corresponding normalized sums of absolute values are
uniformly bounded. For jointly Schwartz \(G(a,h)\) on
\(\mathbb R^k\times\mathbb R^k\), and every \(M\),
\begin{equation}\label{nil:riemann-weighted}
 \sup_h(1+|h|)^M
 \left|v_y\sum_{a\in\mathcal L_y}G(a,h)
                    -\int G(a,h)\,da\right|\le C_{G,M}y.
\end{equation}
\end{lemma}

\begin{proof}
Partition Euclidean space into cells
\(Q_a=a+\prod_j[0,y^{w_j})\).
They have volume \(v_y\) and diameter at most \(\sqrt{k}y\).
For \(x\in Q_a\), the Mean Value Theorem gives
\[
 |G(a,h)-G(x,h)|
     \le \sqrt{k}y\sup_{z\in Q_a}|\nabla_aG(z,h)|
     \le C_{K,M}y(1+|a|)^{-K}(1+|h|)^{-M}.
\]
On each cell \(1+|a|\) and \(1+|x|\) are uniformly comparable.
For \(K>k\), it follows that
\[
 v_y\sum_{a\in\mathcal L_y}(1+|a|)^{-K}
     \le C\int_{\mathbb R^k}(1+|x|)^{-K}\,dx<\infty.
\]
Integrating the cellwise difference and summing proves
Equation~\ref{nil:riemann-weighted}. The same comparison
proves the absolute-value bound and the assertion without \(h\).
\end{proof}

\begin{proposition}[Nonzero discrete twisted projections]\label{nil:projection}
There are constants \(\theta_0>0,C<\infty\), depending only on
\(m\) and the two fixed Schwartz functions, with the following
properties. For \(0<|\theta|<\theta_0\), set
\[
 \epsilon=\operatorname{sign}\theta,\quad
 y=|\theta|^{1/m},\quad
 \nu=y^{mn_{\mathrm f}}=|\theta|^{n_{\mathrm f}},
\]
and define the norm-convergent series
\begin{equation}\label{nil:K}
 K_\theta=\sum_{\gamma\in L}
          \nu J_\epsilon(\Delta_y\gamma)u_\theta(\gamma).
\end{equation}
It is self-adjoint and satisfies, uniformly for both signs,
\begin{align}
 \|K_\theta\|_{\mathrm{coeff},1}&\le C,&
 \|K_\theta\|_{\mathrm{coeff},2}&\asymp\sqrt{\nu},
                                               \label{nil:Knorms}\\
 \|K_\theta^2-K_\theta\|_{\mathrm{coeff},1}&\le Cy,&
 \|K_\theta^2-K_\theta\|_{\mathrm{coeff},2}&\le Cy\sqrt{\nu}.
                                               \label{nil:defects}
\end{align}
Continuous functional calculus in \(C_r^*(L,\theta c)\) gives a
projection \(p_\theta\) with
\begin{equation}\label{nil:projection-bounds}
 \|p_\theta-K_\theta\|\le Cy,\qquad
 0<\operatorname{tr}_\theta(p_\theta)\le C|\theta|^{n_{\mathrm f}}.
\end{equation}
These fields are measurable on the intervals \(0<|\theta|<\theta_0\).
\end{proposition}

\begin{proof}
In quotient coordinates, \(\Delta_y L\) is exactly the
rectangular lattice with coordinate spacings prescribed by the
weights. Its cell volume is \(\nu\), by Lemma~\ref{nil:grading}.
Lemma~\ref{nil:riemann}, applied to \(J_\epsilon\) and
\(|J_\epsilon|^2\), proves Equation~\ref{nil:Knorms}; the
normalized squared sum tends to \(\|J_\epsilon\|_2^2=1\).
Self-adjointness is exact. We have
\[
 J_\epsilon(h)=\overline{J_\epsilon(h^{-1})}
                         \ee(-\epsilon c(h^{-1},h)),
\]
\(\Delta_y(\gamma^{-1})=(\Delta_y\gamma)^{-1}\), and
\[
 \epsilon c(\Delta_y\gamma^{-1},\Delta_y\gamma)
                         =\theta c(\gamma^{-1},\gamma).
\]
Together with the adjoint formula for \(u_\theta\), these
identities prove \(K_\theta^*=K_\theta\) coefficientwise.

For output \(\gamma\), let \(h=\Delta_y\gamma\).
The coefficient of \(K_\theta^2\), divided by \(\nu\), is
\[
 \nu\sum_{a\in\Delta_yL}
   J_\epsilon(a)J_\epsilon(a^{-1}h)
       \ee\bigl(\epsilon c(a,a^{-1}h)\bigr).
\]
Its integral counterpart is \(J_\epsilon(h)\).
The summand is a fixed jointly Schwartz function of \((a,h)\),
as proved in Lemma~\ref{nil:schwartz}. Let \(D_y(h)\) denote
the difference between this sum and \(J_\epsilon(h)\).
Equation~\ref{nil:riemann-weighted} gives
\[
 |D_y(h)|\le C_M y(1+|h|)^{-M}
\]
for every \(M\). The defect coefficient is
\(\nu D_y(\Delta_y\gamma)\). The uniform lattice-sum bound
therefore gives
\[
 \nu\sum_{h\in\Delta_yL}|D_y(h)|\le Cy,\qquad
 \nu^2\sum_{h\in\Delta_yL}|D_y(h)|^2\le Cy^2\nu,
\]
which proves both estimates in Equation~\ref{nil:defects}.

Coefficient \(\ell^1\) controls operator norm, so
\(\|K_\theta^2-K_\theta\|\le Cy\).
For small \(y\), the real spectrum of \(K_\theta\) lies in
\([-1/4,1/4]\cup[3/4,5/4]\), by
\(|\lambda^2-\lambda|\le Cy\) on its spectrum.
Apply one fixed continuous real function equal to zero on
the first interval and one on the second. It gives a
self-adjoint projection \(p_\theta\) in the actual norm-closed
twisted algebra. On these intervals the difference of this
function from \(\lambda\) is \((\lambda^2-\lambda)f(\lambda)\),
where respectively
\[
 f(\lambda)=(1-\lambda)^{-1},\qquad f(\lambda)=-\lambda^{-1}.
\]
Since \(|f|\le2\), functional calculus proves the operator
approximation in Equation~\ref{nil:projection-bounds}.

The upper part is nonempty. If \(p_\theta=0\), then
\(\|K_\theta\|\le1/4\), so
\[
 \|(K_\theta^2-K_\theta)\delta_1\|
 \ge \|K_\theta\delta_1\|-\|K_\theta^2\delta_1\|
 \ge\tfrac34\|K_\theta\delta_1\|.
\]
This contradicts Equations~\ref{nil:Knorms} and
\ref{nil:defects} for small \(y\).
On the two spectral intervals, scalar functional calculus
also gives \(0\le p_\theta\le4K_\theta^2\). Hence
\[
 \operatorname{tr}_\theta(p_\theta)
 \le4\operatorname{tr}_\theta(K_\theta^2)
 =4\|K_\theta\delta_1\|^2\le C\nu.
\]
Strict positivity follows from Lemma~\ref{nil:trace} and the
nonzero projection just established.

Every matrix entry of every finite partial sum defining
\(K_\theta\) is measurable in \(\theta\). Pointwise
coefficient \(\ell^1\) convergence yields strong measurability
on the separable Hilbert space; the common bound yields a
bounded measurable field. Uniform polynomial approximation
of the one chosen continuous spectral function on a common
compact interval gives measurability of \(p_\theta\).
\end{proof}

\begin{lemma}[Compact templates and relative scale changes]\label{nil:templates}
The functions \(J_\epsilon\) in Equation~\ref{nil:K} may be
replaced by smooth compactly supported \(J_\epsilon^\flat\)
at arbitrarily small uniform coefficient \(\ell^1\) cost
for sufficiently small \(|\theta|\).
For fixed such a function, replacing \(\alpha>0\) by \(\beta>0\),
where \(|\log(\alpha/\beta)|\le a\le1\), changes the coefficients
\[
 \alpha^{n_{\mathrm f}}
     J_\epsilon^\flat(\Delta_{\alpha^{1/m}}\gamma)
\]
by \(\ell^1\) amount at most \(Ca\), uniformly for sufficiently
small \(\alpha,\beta\). These coefficient bounds imply operator
bounds when evaluated using any family of unitaries.
\end{lemma}

\begin{proof}
Multiply \(J_\epsilon\) by a smooth cutoff equal to one on a
large coordinate ball. Schwartz decay and the cell comparison
of Lemma~\ref{nil:riemann} make the normalized lattice sum of
the discarded tail arbitrarily small, uniformly as the mesh
goes to zero. For the second assertion set \(r=\alpha/\beta\)
and \(h=\Delta_{\beta^{1/m}}\gamma\). After extracting
\(\beta^{n_{\mathrm f}}\), the difference is
\[
 r^{n_{\mathrm f}}J_\epsilon^\flat(\Delta_{r^{1/m}}h)
                       -J_\epsilon^\flat(h).
\]
For \(e^{-1}\le r\le e\) it is supported on one compact set,
and its supremum is at most \(C|\log r|\) by the Mean Value
Theorem. The volume-normalized number of lattice points there
is uniformly bounded by the cell comparison. Sum to conclude.
\end{proof}

\subsection{A finite magnetic factor and its exact multiplicity}

A single small twist makes the trace small near one endpoint of a
character interval. We need it small near both endpoints, because the
supporting layers will escape in opposite directions there. The finite
representation below shifts the second effective twist from \(t\) to
\(t-\delta\). The product projection then has trace bounded by a
constant times \(N t^{n_{\mathrm f}}(\delta-t)^{n_{\mathrm f}}\).
The prescribed scales will make this small even after the ordinary
matrix multiplicity \(N\) and the geometric weights are included.


The global constant \(h_*=5003\) satisfies \(H_0\le h_*\log_2d\).
After \(m,C_*,C_\bullet\) are fixed, choose
\begin{equation}\label{nil:parameter-order}
 \begin{split}
 &E_0>10(C_*+C_\bullet+1),\\
 &A_*\in\mathbb N,\qquad A_*/m>10E_0h_*,\\
 &\kappa>0,\qquad
        \kappa/m>10(E_0+A_*+C_\bullet+1).
 \end{split}
\end{equation}
Set
\begin{equation}\label{nil:q}
 q=d^{A_*},\qquad \delta=q^{-1},\qquad
 N=q^{n_{\mathrm f}},\qquad I=H_0.
\end{equation}
The numbers \(q,N,I\) are integers. We take \(\ell\) large enough
that \(\delta<\theta_0\) in Proposition~\ref{nil:projection}.

Let \(\mathfrak n_q=\mathfrak n_{\mathbb Z/q\mathbb Z}\), a set of
\(q^{n_{\mathrm f}}=N\) elements. On
\(\mathbb C^N=\ell^2(\mathfrak n_q)\) define a representation of
the second copy \(E_{\mathbb Z}\) by
\begin{equation}\label{nil:R}
 (R(\lambda,b,z)\psi)(x)
   =\ee\bigl(-(z+\lambda(x))/q\bigr)
               \psi\bigl(b^{-1}(1+x)-1\bigr).
\end{equation}
Use any integer representative in the phase; the exponential
is independent of the representative. All arguments on the
right are reduced modulo \(q\). The affine map is a permutation
and the phase has modulus one, so the operator is unitary.
Substitution into Equation~\ref{nil:affine-law}, exactly as for
\(\pi_\epsilon\), proves that \(R\) is a group representation.
Its central character is
\[
 R(0,1,z)=\ee(-\delta z)1_N.
\]
Write \(R_\gamma=R(\sigma(\gamma))\).

Reduction of the integer polynomial laws gives a surjective
homomorphism \(L=\bar E_{\mathbb Z}\to\bar E_{\mathbb Z/q\mathbb Z}\).
Its kernel is
\[
 \Lambda_q=\{(\lambda,b)\in L:
       \lambda_{ij}\equiv(b-1)_{ij}\equiv0\pmod q\}.
\]
It has index \(q^{2n_{\mathrm f}}\) and coset representatives
with every coordinate in \(\{0,\ldots,q-1\}\).
Equation~\ref{nil:R} gives exactly
\begin{equation}\label{nil:congruence-R}
 R_y=1_N\qquad(y\in\Lambda_q).
\end{equation}

The next lemma is a form of twisted absorption
\cite[Lemma~2.4]{EchterhoffWilliams2012}. We track the quotient-identity
block in its conjugacy to retain the exact unnormalized matrix-trace
multiplicity.

\begin{lemma}[Absorption of the finite factor]\label{nil:absorption}
On the second quotient factor, define
\[
 V_t(\gamma)=u_t(\gamma)\otimes R_\gamma.
\]
Its multiplier is \(\ee((t-\delta)c)\), and it is unitarily
equivalent to \(N\) copies of \(u_{t-\delta}\).
The equivalence fixes the quotient-identity fiber, preserving
the identity trace with its unnormalized matrix sum.
Thus every coefficient sum, and every continuous spectral
function of it, has the same operator norm as its evaluation
in \(u_{t-\delta}\); the identity trace of a projection is
multiplied by \(N\).
\end{lemma}

\begin{proof}
The section law and the central character of \(R\) give
\[
 R_\gamma R_\eta
       =\ee(-\delta c(\gamma,\eta))R_{\gamma\eta}.
\]
Combining this with the multiplier of \(u_t\) proves the
multiplier assertion. Define a diagonal unitary by
\[
 W_t(\delta_\eta\otimes v)=\delta_\eta\otimes R_\eta v.
\]
Its formula is actually independent of \(t\). Then
\begin{align*}
 V_t(\gamma)W_t(\delta_\eta\otimes v)
  &=\ee(t c(\gamma,\eta))
                 \delta_{\gamma\eta}\otimes R_\gamma R_\eta v\\
  &=\ee((t-\delta)c(\gamma,\eta))
                 \delta_{\gamma\eta}\otimes R_{\gamma\eta}v\\
  &=W_t(u_{t-\delta}(\gamma)\otimes1_N)
                 (\delta_\eta\otimes v).
\end{align*}
So \(W_t^*V_t(\gamma)W_t=u_{t-\delta}(\gamma)\otimes1_N\).
As \(R_1=1_N\), \(W_t\) fixes \(\delta_1\otimes\mathbb C^N\)
pointwise, and conjugacy preserves the entire identity matrix
block. All conclusions follow first for finite sums and then
by norm limits and continuous functional calculus.
\end{proof}

In the full Fourier quotient space, amplified by \(\mathbb C^N\),
the first factor \(u_t\) and the second \(V_t\) commute. They are
equivariant for the original right quotient actions tensored
with \(1_N\).
For \(0<t<\delta\), both \(t\) and \(t-\delta\) fall in
Proposition~\ref{nil:projection}. Define
\begin{equation}\label{nil:pt}
 p(t)=p_t\otimes p_{t-\delta}[V_t].
\end{equation}
The second-factor notation means coefficient evaluation and
continuous functional calculus through Lemma~\ref{nil:absorption}.
Set \(p(t)=0\) outside this interval.
It is a measurable bounded field of equivariant projections.
Its identity trace \(s_p(t)=\operatorname{Tr}_{\mathrm{id},N}(p(t))\)
satisfies
\begin{equation}\label{nil:two-factor-trace}
 \begin{split}
 0<s_p(t)
 &=N\,\operatorname{tr}_t(p_t)
               \operatorname{tr}_{t-\delta}(p_{t-\delta})\\
 &\le C N t^{n_{\mathrm f}}(\delta-t)^{n_{\mathrm f}}
                  \qquad(0<t<\delta).
 \end{split}
\end{equation}
Factor the identity diagonals of the tensor factors and use
the identity-block assertion of Lemma~\ref{nil:absorption}
for the equality. In subsequent coefficient formulas, an
additional shared central power \(z_0^n\) acts by \(\ee(nt)1_N\).
The matrix \(R_\gamma\) is attached to a second-factor section
term as in \(V_t\); these prescriptions fix all phases.

\subsection{The layer window and its weighted trace}

\begin{lemma}[Smooth layer windows]\label{nil:window}
There is a smooth nonnegative \(h\), supported in \((-1,1)\)
and positive on \([-0.7,0.7]\), with
\(\sum_{i\in\mathbb Z}h(v-i)^2=1\) for every \(v\in\mathbb R\).
For \(0<t<\delta\), set
\begin{equation}\label{nil:h}
 b(t)=\kappa^{-1}\log_2\frac{t}{\delta-t},\qquad
 h_i(t)=h(b(t)-i),
\end{equation}
and extend \(h_i\) by zero to the rest of the circle.
Each extension is smooth. The active set
\(J(t)=\{i:h_i(t)\ne0\}\) has one or two elements, adjacent
if two, and \(\sum_i h_i(t)^2=1\) on the open interval.
With \(w_i=q^{-1}2^{-\kappa|i|}\), its support length and
derivatives satisfy
\[
 |\operatorname{supp}h_i|\le C_\kappa w_i,\qquad
 \|\partial_t^j h_i\|_\infty\le C_{\kappa,j}w_i^{-j}.
\]
On an active index \(i\), uniformly in \(\ell,i,t\),
\begin{align}
 \log_2(t^{-1})
 &=\log_2q+\kappa\max(-i,0)+O_\kappa(1),\label{nil:log-t}\\
 \log_2((\delta-t)^{-1})
 &=\log_2q+\kappa\max(i,0)+O_\kappa(1),\label{nil:log-tau}\\
 t(\delta-t)&\le 2^\kappa q^{-2}2^{-\kappa|i|}.
                                                  \label{nil:product}
\end{align}
\end{lemma}

\begin{proof}
Choose a smooth nonnegative bump \(\beta\), supported in
\((-0.9,0.9)\) and positive on \([-0.7,0.7]\), and put
\[
 h(v)=\beta(v)\left(\sum_{j\in\mathbb Z}\beta(v-j)^2\right)^{-1/2}.
\]
The denominator is positive, smooth, and one-periodic:
some translate lies in \([-1/2,1/2]\).
This proves the square-sum property. An interval of length
\(1.8<2\) contains at most two integers, necessarily adjacent.
This proves the assertion about active indices.
For each fixed \(i\), its support in \(b\)-coordinates lies
in a compact interval strictly inside the inverse image
of \((0,\delta)\); its \(t\)-support is therefore a compact
subset of \((0,\delta)\), so extension by zero is smooth.

The inverse map and its derivative are
\[
 t(b)=\delta\frac{2^{\kappa b}}{1+2^{\kappa b}},\qquad
 \delta-t(b)=\frac{\delta}{1+2^{\kappa b}},\qquad
 t'(b)=\kappa(\log2)\frac{t(b)(\delta-t(b))}{\delta}.
\]
For \(|b-i|\le0.9\), \(t'(b)\) is bounded above and below by
positive constants depending on \(\kappa\) times \(w_i\).
Every higher \(b\)-derivative is bounded by a constant times
\(w_i\), by differentiating the displayed rational function
and separating \(b\ge0\) from \(b\le0\).
The inverse-function differentiation formulas and the fixed
derivatives of \(h\) give the stated derivative bounds.
Integration of the bound on \(t'\) proves the support bound.

The exact logarithmic identities are
\[
 \log_2(t(b)^{-1})=\log_2q+\log_2(1+2^{-\kappa b}),\quad
 \log_2((\delta-t(b))^{-1})=\log_2q+\log_2(1+2^{\kappa b}).
\]
For real \(x\), \(\log_2(1+2^x)-\max(x,0)\) lies in \([0,1]\).
Using \(|b-i|<1\) proves Equations~\ref{nil:log-t} and
\ref{nil:log-tau}. Finally
\[
 \frac{2^{\kappa b}}{(1+2^{\kappa b})^2}\le2^{-\kappa|b|}
\]
and \(|b-i|<1\) prove Equation~\ref{nil:product}.
\end{proof}

\begin{proposition}[Uniform weighted trace bound]\label{nil:weighted}
For \(\mathfrak w_i=2^{C_\bullet H_i}\), define
\[
 \mu_\ell=\sup_{\substack{0<t<\delta\\i\in J(t)}}
                          \mathfrak w_i\,s_p(t).
\]
Then \(\mu_\ell\to0\) as \(\ell\to\infty\). More explicitly,
\begin{equation}\label{nil:weighted-estimate}
 \mathfrak w_i N t^{n_{\mathrm f}}(\delta-t)^{n_{\mathrm f}}
 \le C d^{-(A_*n_{\mathrm f}-C_\bullet h_*)}
          2^{-(\kappa n_{\mathrm f}-C_\bullet)|i|}
                  \qquad(i\in J(t)).
\end{equation}
The two exponents subtracted on the right are positive.
\end{proposition}

\begin{proof}
Using \(H_i=H_0+|i|\), \(H_0\le h_*\log_2d\),
\(N=q^{n_{\mathrm f}}\), and Equation~\ref{nil:product},
\begin{align*}
 2^{C_\bullet H_i}N\bigl(t(\delta-t)\bigr)^{n_{\mathrm f}}
 &\le 2^{\kappa n_{\mathrm f}}
       2^{C_\bullet H_0}q^{-n_{\mathrm f}}
       2^{-(\kappa n_{\mathrm f}-C_\bullet)|i|}\\
 &\le 2^{\kappa n_{\mathrm f}}
       d^{-(A_*n_{\mathrm f}-C_\bullet h_*)}
       2^{-(\kappa n_{\mathrm f}-C_\bullet)|i|}.
\end{align*}
Equation~\ref{nil:parameter-order}, with \(m\ge2\) and
\(n_{\mathrm f}\ge1\), implies
\(A_*n_{\mathrm f}>C_\bullet h_*\) and
\(\kappa n_{\mathrm f}>C_\bullet\).
Take the supremum and apply Equation~\ref{nil:two-factor-trace}.
\end{proof}

The dependency order for subsequent arguments is as follows.
Geometry fixes \(C_*,C_\bullet\); Equation~\ref{nil:choice-m}
then fixes \(m\), and Lemma~\ref{nil:schwartz} fixes two functions.
Equation~\ref{nil:parameter-order} fixes \(E_0,A_*,\kappa\).
All accuracies, compact truncations, relative logarithmic mesh
sizes, and finite norm-test exponents may now be chosen as
fixed constants. Only then is \(\ell\) increased. In particular
the trace estimate requires no accuracy varying with the layer,
and the word detector does not depend on subsequent moment orders.
